\documentclass[11pt,a4paper]{article}

% Fonts
\usepackage[utf8]{inputenc}
\usepackage[T1]{fontenc}
\usepackage{textcomp}
\usepackage{lmodern}

% Layout
\usepackage[margin=2.5cm]{geometry}
\usepackage{setspace}
\setlength{\parindent}{0pt}
\setlength{\parskip}{0.7em plus 0.2em minus 0.1em}
\onehalfspacing

% Language
\usepackage{babel}
\babelprovide[main,import]{british}

% Bibliography (APA 7th edition)
\usepackage{csquotes}
\usepackage[backend=biber,style=apa,sorting=nyt,natbib=true]{biblatex}
\DeclareLanguageMapping{british}{british-apa}
\addbibresource{bibliography.bib}

% Mathematics
\usepackage{amsmath,amssymb}

% Tables and figures
\usepackage{booktabs}
\usepackage{longtable}
\usepackage{graphicx}
\usepackage{float}
\usepackage{array}
\usepackage{multirow}
\usepackage{tabularx}
\usepackage{pdflscape}
\usepackage{enumitem}

\newcolumntype{L}[1]{>{\raggedright\arraybackslash}p{#1}}
\newcolumntype{C}[1]{>{\centering\arraybackslash}p{#1}}

% Hyperlinks
\usepackage{xcolor}
\usepackage[colorlinks=true,linkcolor=blue!60!black,citecolor=blue!60!black,urlcolor=blue!60!black]{hyperref}

% Captions
\usepackage[font=small,labelfont=bf]{caption}

% Section formatting
\usepackage{titlesec}
\titleformat{\section}{\large\bfseries}{\thesection}{1em}{}
\titleformat{\subsection}{\normalsize\bfseries}{\thesubsection}{1em}{}
\titleformat{\subsubsection}{\normalsize\itshape}{\thesubsubsection}{1em}{}

% TOC: widen number columns to prevent A.10, A.1.10 etc. from overlapping titles
\makeatletter
\renewcommand{\l@section}{\@dottedtocline{1}{0em}{2.8em}}
\renewcommand{\l@subsection}{\@dottedtocline{2}{2.8em}{3.6em}}
\renewcommand{\l@subsubsection}{\@dottedtocline{3}{6.4em}{4.2em}}
\makeatother

% Custom commands
\newcommand{\pp}{\text{pp}}
\newcommand{\EUR}{\texteuro}
\newcommand{\code}[1]{\texttt{#1}}

% Renumber sections as Appendix letters for standalone compilation
\renewcommand{\thesection}{A.\arabic{section}}
\renewcommand{\thetable}{A.\arabic{table}}
\renewcommand{\thefigure}{A.\arabic{figure}}

% Title
\title{\textbf{Appendix 1:\\  Data sources for Global Minimum Tax thesis}\\[0.5em]
\large MSc Economic Policy Thesis\\Utrecht University School of Economics}
\author{Maurits de Haan\\[0.3em]
\small Supervisor: dr.\ Peter Gerbrands}
\date{March 2026}

% COMPILE ORDER: pdflatex → biber → pdflatex → pdflatex
% (or: latexmk -pdf -bibtex cbcr_data_annex)

\begin{document}

\maketitle

\newpage
\tableofcontents
\newpage

\section{OECD CbCR data}

Under the OECD/G20 Base Erosion and Profit Shifting (BEPS) Action~13, large multinational enterprise (MNE) groups with consolidated revenue exceeding EUR~750~million are required to file Country-by-Country Reports (CbCRs) detailing their economic activity and tax payments in each jurisdiction where they operate.  The OECD aggregates and anonymises these reports into jurisdiction-level statistics published in the \emph{OECD Corporate Tax Statistics} (CTS) database.

The CbCR statistics are uniquely valuable because they capture tax-planning behaviour at the MNE level across \emph{all} jurisdictions, including low-tax locations typically absent from firm-level datasets.  Key measures include pre-tax profits, taxes accrued and paid, employees, total revenue, and assets, all disaggregated by \emph{parent jurisdiction} (REF\_AREA) and \emph{partner jurisdiction} (COUNTERPART\_AREA).

\subsection{Data architecture and SDMX structure}\label{sec:architecture}

The OECD publishes Country-by-Country Reporting statistics through its SDMX REST API under the endpoint \url{https://sdmx.oecd.org/public/rest/data}. All six CbCR tables share an identical column structure and differ only in the combinations of dimension values. Each row of each file represents one observation uniquely identified by the combination of its dimension values; the quantity of interest is stored in the \code{OBS\_VALUE} column.

The data are structured according to the Statistical Data and Metadata eXchange (SDMX 2.1) standard. Every dataset is identified by a \emph{dataflow} code of the form \code{OECD.CTP.TPS,DSD\_CBCR@DF\_CBCR[I--VI],1.1}. The tables cover fiscal years 2016-2022.

\subsection{Common column structure}

All six files share the following fifteen columns:

\begin{table}[H]
\centering
\caption{Column structure common to all six CbCR SDMX files}
\label{tab:columns}
\begin{tabular}{@{}lp{9.5cm}@{}}
\toprule
Column name & Description \\
\midrule
\code{DATAFLOW} & Identifier of the SDMX dataflow and version \\
\code{REF\_AREA} & ISO 3166-1 alpha-3 code of the \emph{reporting (parent) jurisdiction}: the country whose tax authority collected and submitted the CbCR data \\
\code{COUNTERPART\_AREA} & ISO 3166-1 alpha-3 code of the \emph{affiliate jurisdiction}: the country where the constituent entity is resident; \code{\_Z} if not applicable; regional aggregates such as \code{WXD} (world excl.\ domestic), \code{W} (world total), and continent codes also appear \\
\code{FREQ} & Frequency; always \code{A} (annual) \\
\code{MEASURE} & The variable being measured (see Table~\ref{tab:measures}) \\
\code{UNIT\_MEASURE} & Unit of the observation (see Table~\ref{tab:units}) \\
\code{STATISTICAL\_OPERATION} & Statistical transformation applied, if any (see Table~\ref{tab:statops}) \\
\code{AGGREGATION\_TYPE} & Unit of observation and aggregation level (see Table~\ref{tab:aggtypes}) \\
\code{PROFIT\_GROUPING} & Panel filter (all MNEs, profit-makers only, all incl.\ loss-makers) \\
\code{EFFECTIVE\_TAX\_RATE} & ETR bracket classification; \code{\_Z} if not disaggregated by ETR \\
\code{MNE\_SIZE} & Revenue-based size class of the MNE group; \code{\_Z} if not disaggregated \\
\code{TIME\_PERIOD} & Fiscal year of the reported data \\
\code{OBS\_VALUE} & The observation value in the unit specified by \code{UNIT\_MEASURE}; monetary values are denominated in USD \\
\code{DECIMALS} & Number of decimal places used to express the value \\
\code{OBS\_STATUS} & \code{A} = actual/normal value; \code{Q} = provisional or preliminary value \\
\bottomrule
\end{tabular}
\end{table}

\subsection{Missing-data convention}

In the analysis to find how the missing-data was handled in the dataset, I found out that the difference lies in how \code{OBS\_VALUE} looks:

\begin{itemize}
    \item \code{OBS\_VALUE} = \textit{numerical value} (including zero) combined with
          \code{OBS\_STATUS = A}: the observation is published and available. A
          zero value means no CbCR activity was reported for that dimension combination,
          which is substantively meaningful (i.e.\ the absence of a bilateral
          relationship is a real fact about the data).
    \item \code{OBS\_VALUE} = \textit{missing} (\texttt{NaN}) combined with
          \code{OBS\_STATUS = Q}: the observation exists but has been suppressed, is
          not yet finalised, or has been withheld for confidentiality reasons.
\end{itemize}

This convention means that the true missing rate (NaN entries) is substantially lower than one might expect given the SDMX ``sparse'' storage format. Across the six tables, the NaN rate ranges from 1.1\% (Table~V) to 6.5\% (Table~IV); see Table~\ref{tab:crosssummary} for details.

\subsection{Dimension code reference}\label{sec:dimensions}

Table~\ref{tab:measures} lists all \code{MEASURE} codes that appear across the six CbCR tables, together with their substantive interpretation and the unit in which they are reported. Not all measures appear in all tables; the column \emph{Tables} indicates where each measure is available.

\begin{longtable}{@{}lp{6cm}lp{2.5cm}@{}}
\caption{MEASURE codes in the OECD CbCR database}\label{tab:measures}\\
\toprule
Code & Description & Unit & Tables \\
\midrule
\endfirsthead
\multicolumn{4}{c}{\textit{Table~\ref{tab:measures} continued}}\\
\toprule
Code & Description & Unit & Tables \\
\midrule
\endhead
\midrule
\multicolumn{4}{r}{\textit{Continued on next page}}\\
\endfoot
\bottomrule
\endlastfoot
\multicolumn{4}{l}{\textit{Financial aggregates}}\\
\code{PROFIT}         & Profit before income tax. The primary measure of taxable profit used in CbCR; susceptible to profit shifting and earnings management & USD & I--IV, VI \\
\code{PROFIT\_ADJ}    & Adjusted profit before tax (corrected for intra-group dividends and equity income to reduce double-counting) & USD & I--IV \\
\code{TAX\_ACCRUED}   & Income tax accrued (current year charge); the numerator of the CbCR effective tax rate & USD & I--IV, VI \\
\code{TAX\_PAID}      & Income tax paid on a cash basis; may differ from accrued due to timing of payments & USD & I--IV, VI \\
\code{TOT\_REV}       & Total revenue: the sum of unrelated-party and related-party revenues & USD & I--IV, VI \\
\code{UPR}            & Unrelated-party revenue; excludes intra-group transactions & USD & I--VI \\
\code{RPR}            & Related-party revenue; intra-group revenues between constituent entities & USD & I--IV, VI \\
\code{ASSETS}         & Tangible assets, excluding cash and cash equivalents; the denominator of the substance-based income exclusion (payroll component) & USD & I--VI \\
\code{EARNINGS}       & Total payroll costs (compensation of employees); the second component of the SBIE denominator & USD & I--IV, VI \\
\code{STATED\_CAPITAL}& Stated capital and accumulated retained earnings & USD & I--IV, VI \\
\addlinespace
\multicolumn{4}{l}{\textit{Count variables}}\\
\code{EMPLOYEES}      & Number of employees (full-time equivalent); an alternative substance proxy & Persons & I--VI \\
\code{ENTITIES\_COUNT}& Number of constituent entities reporting in that jurisdiction & Entities & I--IV \\
\code{CBCR\_COUNT}    & Number of CbCR filings (MNE groups) contributing to the aggregate & MNE groups & I--IV, VI \\
\code{SUBGROUPS\_COUNT}& Number of MNE subgroups & Subgroups & I--IV, VI \\
\addlinespace
\multicolumn{4}{l}{\textit{Derived tax rate}}\\
\code{ETR}            & Effective tax rate computed as \code{TAX\_ACCRUED} / \code{PROFIT} at the jurisdiction level; expressed as a percentage of profit before tax & \% of PBT & I \\
\addlinespace
\multicolumn{4}{l}{\textit{Activity-share variables (all in Table I, III, IV, VI)}}\\
\code{ACT\_ADMIN}     & Share of entities engaged in administrative or management services & \% / est.\ & I, III, IV, VI \\
\code{ACT\_DOR}       & Share of entities that are dormant & \% / est.\ & I, III, IV, VI \\
\code{ACT\_FIN}       & Share of entities engaged in internal finance or insurance & \% / est.\ & I, III, IV, VI \\
\code{ACT\_HOLDING}   & Share of entities classified as holding companies & \% / est.\ & I, III, IV, VI \\
\code{ACT\_IGF}       & Share of entities engaged in intra-group financing & \% / est.\ & I, III, IV, VI \\
\code{ACT\_INS}       & Share of entities engaged in insurance & \% / est.\ & I, III, IV, VI \\
\code{ACT\_IP}        & Share of entities holding intellectual property & \% / est.\ & I, III, IV, VI \\
\code{ACT\_MAN}       & Share of entities engaged in manufacturing or production & \% / est.\ & I, III, IV, VI \\
\code{ACT\_OTHER}     & Share of entities engaged in activities not classified elsewhere & \% / est.\ & I, III, IV, VI \\
\code{ACT\_PUR}       & Share of entities engaged in purchasing or procurement & \% / est.\ & I, III, IV, VI \\
\code{ACT\_RD}        & Share of entities engaged in research and development & \% / est.\ & I, III, IV, VI \\
\code{ACT\_SALES}     & Share of entities engaged in sales, marketing or distribution & \% / est.\ & I, III, IV, VI \\
\code{ACT\_SER}       & Share of entities providing services (excluding IP and R\&D) & \% / est.\ & I, III, IV, VI \\
\end{longtable}

\subsection{Unit codes (\texttt{UNIT\_MEASURE})}

\begin{table}[H]
\centering
\caption{UNIT\_MEASURE codes}
\label{tab:units}
\begin{tabular}{@{}llp{8cm}@{}}
\toprule
Code & Abbreviation & Description \\
\midrule
\code{USD}            & USD & US dollars (nominal); monetary variables are expressed in full USD (not thousands or millions) \\
\code{PS}             & Persons & Number of individuals (for EMPLOYEES) \\
\code{EST}            & Entities & Number of entities, or activity shares expressed as estimated percentages \\
\code{ENT\_MULT}      & Multiples & Number of entity-multiples (for count variables in some aggregations) \\
\code{PT\_PRFT\_BF\_TAX} & \% PBT & Percentage of profit before tax (used for the ETR measure) \\
\bottomrule
\end{tabular}
\end{table}

\subsection{Statistical operation codes (\texttt{STATISTICAL\_OPERATION})}

\begin{table}[H]
\centering
\caption{STATISTICAL\_OPERATION codes}
\label{tab:statops}
\begin{tabular}{@{}llp{9.5cm}@{}}
\toprule
Code & Label & Description \\
\midrule
\code{\_Z} & Aggregate & Simple aggregate total; no distributional transformation applied \\
\code{P5}  & 5th pct.\ & 5th percentile of the distribution of MNE groups \\
\code{P10} & 10th pct.\ & 10th percentile \\
\code{P25} & 25th pct.\ & Lower quartile (25th percentile) \\
\code{P50} & Median & Median (50th percentile) \\
\code{P75} & 75th pct.\ & Upper quartile (75th percentile) \\
\code{P90} & 90th pct.\ & 90th percentile \\
\code{P95} & 95th pct.\ & 95th percentile \\
\code{IQMEAN} & IQ mean & Interquartile mean: the arithmetic mean computed after excluding the bottom and top 25\% of observations; a robust estimator less sensitive to outlier MNEs \\
\bottomrule
\end{tabular}
\end{table}

\subsection{Aggregation type codes (\texttt{AGGREGATION\_TYPE})}

The aggregation type defines the unit of observation at which the CbCR data are
compiled. 

\begin{table}[H]
\centering
\caption{AGGREGATION\_TYPE codes}
\label{tab:aggtypes}
\begin{tabular}{@{}llp{8cm}@{}}
\toprule
Code & Label & Description \\
\midrule
\code{TOTAL}   & Jurisdiction total & Financial variables summed across all MNE groups reporting to the parent jurisdiction's tax authority; the standard aggregate used in \citet{barake2022} \\
\code{TRMNE}   & MNE-group level & Variables computed at the level of the individual MNE group and then cross-tabulated by the reporting (parent) jurisdiction; used in Table III to distribute entities across ETR brackets by MNE group \\
\code{TRSUB}   & Subsidiary level & Variables computed at the level of the individual constituent entity (subsidiary) and aggregated by its location jurisdiction; used in Table IV to identify where low-taxed profit is booked \\
\code{MNESIZE} & Size class & Variables aggregated within each of five revenue-based size classes of MNE groups; used exclusively in Table VI \\
\bottomrule
\end{tabular}
\end{table}

\subsection{Effective tax rate bracket codes (\texttt{EFFECTIVE\_TAX\_RATE})}

The ETR brackets are the central dimension of Tables III and IV. Table~\ref{tab:etr}
documents all codes, with particular attention to the brackets most relevant for
GMT revenue simulation.

\begin{table}[H]
\centering
\caption{EFFECTIVE\_TAX\_RATE bracket codes}
\label{tab:etr}
\begin{tabular}{@{}lp{11cm}@{}}
\toprule
Code & Description \\
\midrule
\code{\_Z}                         & Not applicable; the observation is not disaggregated by ETR \\
\code{NEGATIVE\_PROFIT}            & Entities (or groups) with negative profit before tax (loss-makers) \\
\code{NEGATIVE\_PROFIT\_TAX\_ACCRUED} & Entities with both negative profit and negative tax accrued \\
\code{NEGATIVE\_TAX\_ACCRUED}      & Entities with positive profit but negative tax accrued (i.e.\ a tax credit or refund) \\
\code{RATE\_0-10}                  & ETR strictly below 10\%; the most aggressive avoidance bracket \\
\code{RATE\_0-15}                  & ETR below 15\%; the \textbf{GMT-critical bracket}: all profit in this range triggers the Pillar Two top-up tax \\
\code{RATE\_10}                    & ETR between 10\% and 15\% (paid-tax basis) \\
\code{RATE\_10\_ACCRUED}           & ETR between 10\% and 15\% (accrued-tax basis) \\
\code{RATE\_15}                    & ETR between 15\% and 20\% \\
\code{RATE\_20}                    & ETR between 20\% and 30\% (summary bracket) \\
\code{RATE\_10-20}                 & ETR between 10\% and 20\% \\
\code{RATE\_20-30}                 & ETR between 20\% and 30\% \\
\code{RATE\_30}                    & ETR above 30\% \\
\bottomrule
\end{tabular}
\end{table}

\subsection{Profit grouping codes (\texttt{PROFIT\_GROUPING})}

\begin{table}[H]
\centering
\caption{PROFIT\_GROUPING codes}
\label{tab:panelgroups}
\begin{tabular}{@{}llp{9cm}@{}}
\toprule
Code & Label & Description \\
\midrule
\code{PANELAI}  & Panel A-I & Restricts the sample to MNE groups (or constituent entities) with positive profit before tax. Used to compute the ETR and revenue estimates. Revenue estimates based solely on Panel A-I understate potential GMT revenue because loss-making entities may still trigger top-up taxes at the jurisdiction level \\
\code{PANELAII} & Panel A-II & Includes all MNE groups, including those with aggregate losses. Preferred for total activity coverage but not suitable as the ETR denominator \\
\code{\_T}       & All MNEs & Combined total of both panels; equivalent to Panel A-II in most contexts \\
\code{\_Z}       & N/A & Panel grouping not applicable for this observation \\
\bottomrule
\end{tabular}
\end{table}

\subsection{MNE size class codes (\texttt{MNE\_SIZE})}

The size classes are defined by consolidated total revenue (not profit) of the
ultimate parent entity's MNE group. The EUR~750~million CbCR threshold means all
reporting MNEs are, by construction, at the high end of the revenue distribution.

\begin{table}[H]
\centering
\caption{MNE\_SIZE codes (Table VI only)}
\label{tab:mnesize}
\begin{tabular}{@{}llp{8cm}@{}}
\toprule
Code & Revenue range & Description \\
\midrule
\code{SIZE\_LESS\_2\_5}   & $<$ USD 2.5 bn & Smallest MNE groups above the CbCR threshold (EUR 750M $\approx$ USD 800M) \\
\code{SIZE\_2\_5-5}       & USD 2.5--5 bn & Medium-small MNE groups \\
\code{SIZE\_5-10}         & USD 5--10 bn & Medium MNE groups \\
\code{SIZE\_GREATER\_10}  & $>$ USD 10 bn & Largest MNE groups; account for the majority of global shifted profits \\
\code{SIZE\_LESS\_5}      & $<$ USD 5 bn & Summary: combines \code{SIZE\_LESS\_2\_5} and \code{SIZE\_2\_5-5} \\
\code{SIZE\_GREATER\_5}   & $>$ USD 5 bn & Summary: combines \code{SIZE\_5-10} and \code{SIZE\_GREATER\_10} \\
\bottomrule
\end{tabular}
\end{table}

\subsection{Cross-Table Overview}\label{sec:crossoverview}

Table~\ref{tab:crosssummary} summarises the six CbCR tables along the key
structural dimensions. The database as a whole contains approximately 2.47 million
rows across all tables. The total file size is approximately 330\,MB (uncompressed
CSV format).

\begin{landscape}
\begin{table}[H]
\centering
\caption{Cross-table summary of OECD CbCR Tables I--VI}
\label{tab:crosssummary}
\footnotesize
\begin{tabular}{@{}lrrrrrrr@{}}
\toprule
 & \textbf{Table I} & \textbf{Table II} & \textbf{Table III} & \textbf{Table IV} & \textbf{Table V} & \textbf{Table VI} \\
 & Agg.\ Totals & IQ Means & ETR (MNE) & ETR (Sub.) & Size Dist.\ & Size Agg.\ \\
\midrule
Rows (total) & 1,124,211 & 437,129 & 699,246 & 101,399 & 4,195 & 105,753 \\
Positive obs (\%) & 71.5 & 84.6 & 68.3 & 77.3 & 98.4 & 72.1 \\
Zero obs (\%) & 21.7 & 7.6 & 24.9 & 12.1 & 0.4 & 22.9 \\
NaN / missing (\%) & 3.0 & 1.9 & 3.6 & 6.5 & 1.1 & 2.4 \\
\addlinespace
Fiscal years covered & 2016--2022 & 2017--2022 & 2017--2022 & 2016--2022 & 2016--2022 & 2022 only \\
Parent jurisdictions & 58 & 46 & 46 & 49 & 40 & 39 \\
Affiliate jurisdictions & 240 & 238 & 238 & 154 & --- & 215 \\
Number of MEASURE codes & 28 & 14 & 27 & 27 & 3 & 26 \\
\addlinespace
Aggregation type & \code{TOTAL} & \code{TOTAL} & \code{TRMNE} & \code{TRSUB} & \code{TRMNE} & \code{MNESIZE} \\
ETR brackets & No & No & Yes (12) & Yes (11) & No & No \\
Profit panels & A-I, A-II & A-I, A-II & --- & --- & --- & --- \\
Statistical operations & Pct., Agg & IQmean & Agg & Agg & Pctiles & Agg \\
MNE size classes & No & No & No & No & No & Yes (6) \\
\bottomrule
\end{tabular}
\end{table}
\end{landscape}

\subsection{Year-by-year expansion of coverage}

The CbCR data shwos an expansion of reporting coverage over time, as more jurisdictions adopted the BEPS Action 13 CbCR requirement. Table~\ref{tab:yearcoverage} documents the number of parent (reporting) jurisdictions and the count of positive observations in Table~I for each fiscal year.

\begin{table}[H]
\centering
\caption{Year-by-year coverage of Table~I (Aggregate Totals)}
\label{tab:yearcoverage}
\begin{tabular}{@{}lrrrrrr@{}}
\toprule
Year & Parent & Total  & Positive & Zero & NaN & Positive \\
     & jurisdictions & rows & obs & obs & obs & rate (\%) \\
\midrule
2016 & 26 &  63,735 &  51,553 &  12,182 &      0 & 80.9 \\
2017 & 38 & 132,477 &  99,780 &  32,697 &      0 & 75.3 \\
2018 & 46 & 153,671 & 116,698 &  36,973 &      0 & 75.9 \\
2019 & 50 & 166,717 & 125,913 &  40,804 &      0 & 75.5 \\
2020 & 52 & 171,924 & 128,725 &  43,199 &      0 & 74.9 \\
2021 & 52 & 236,288 & 137,721 &  68,149 & 30,418 & 58.3 \\
2022 & 54 & 199,399 & 143,140 &  52,749 &  3,510 & 71.8 \\
\bottomrule
\end{tabular}
\end{table}

The most recent release (fiscal year 2022) reports data from 54 parent jurisdictions, representing approximately 98\% of all CbCR filings globally \citep{oecd2025corptaxstats}. The 2016 release (the first year) covered only 26 jurisdictions, and the dataset grew substantially in 2017-2018 as BEPS Action 13 legislation was enacted across OECD and G20 members. The 2021 vintage shows a spike in NaN observations (30,418 rows), which largely disappeared in
2022; this likely reflects a transitional period in OECD data processing after the
redesign of the CbCR publication framework.\footnote{The 2021 vintage was the
first published under the new disaggregated SDMX format; some jurisdictions
submitted incomplete data during this transition.}

\subsection{Table I: Aggregate Totals by Jurisdiction}\label{sec:table1}

\paragraph{Purpose and structure}
Table~I is the primary CbCR table and the starting point for all revenue simulations. It reports jurisdiction-level financial aggregates, broken down by parent jurisdiction (REF\_AREA), affiliate jurisdiction (COUNTERPART\_AREA), profit panel (PANELAI / PANELAII), and statistical operation (aggregate, P5, P25, P50, P75, P95). The combination of parent and affiliate jurisdiction creates a bilateral dataset analogous to a trade matrix: each cell represents the total activity of all MNEs headquartered in country $i$ that have constituent entities in country $j$.

\paragraph{Key features}
With 1,124,211 rows and 28 MEASURE codes, Table~I is the largest of the six tables. It is the only table that contains the \code{ETR} measure (effective tax rate, expressed as a percentage of profit before tax) and the only table that provides full distributional percentiles (P5 through P95) for all measures. The \code{STATISTICAL\_OPERATION = \_Z} rows provide the aggregate totals used in \citet{barake2022}; the percentile rows enable the within-jurisdiction ETR heterogeneity analysis following \citet{hugger2025}.

\paragraph{Coverage}
In 2022, 54 parent jurisdictions report data, with 143,140 positive observations. The main bilateral region aggregates available as \code{COUNTERPART\_AREA} values include: individual ISO country codes (154 countries), \code{WXD} (world excluding domestic), \code{W} (world total), \code{STLS} (stateless entities), and several regional aggregates (\code{A}, \code{A\_O}, \code{E}, \code{E\_O}, \code{F}, \code{F\_O}, \code{S}, \code{S\_O} representing Americas, Europe, Africa, and similar groupings excluding the parent's own region).

\paragraph{Limitations}
The aggregate totals are vulnerable to the double-counting problem identified by \citet{blouin2023}: when a German parent owns a Dutch holding company which owns an Irish operating subsidiary, the Irish subsidiary's profits may be attributed to both Germany (through the CbCR filing) and the Netherlands (as an intermediate parent). The \code{PROFIT\_ADJ} measure is the OECD's attempt to correct for intra-company dividends, but it is available only for 2021--2022 and for a subset of parent jurisdictions.

\subsection{Table II: Interquartile Means by Jurisdiction}\label{sec:table2}

\paragraph{Purpose and structure}
Table~II provides the interquartile mean (IQM) of financial variables: the arithmetic mean computed after excluding the bottom and top quartiles of MNE groups within each parent-affiliate jurisdiction pair. This is a robust central tendency estimator, less sensitive to the extremely large MNE groups (e.g.\ Apple, ExxonMobil) that can dominate simple averages.

\paragraph{Key features}
Table~II contains 14 MEASURE codes, all financial aggregates from Table~I except activity shares and the ETR measure itself. It has a higher positive observation rate (84.6\%) than Table~I because the IQM is only computed where there are at least four MNE groups, reducing the number of zero cells. The data starts from 2017 (one year later than Table~I) because computing a robust mean requires a minimum number of observations.

\paragraph{Use for methodology}
The IQM is useful for characterising the ``typical'' MNE group in a given jurisdiction pair, as distinct from the aggregate total driven by outlier firms. For the SBIE computation, the IQM of EARNINGS and ASSETS per MNE group provides a cross-check on the aggregate totals in Table~I. 

\subsection{Table III: Entity Counts by ETR Bracket (MNE Level)}\label{sec:table3}

\paragraph{Purpose and structure}
Table~III disaggregates constituent entities by effective tax rate bracket at the \emph{MNE group} level (AGGREGATION\_TYPE = TRMNE). This means that an entity is assigned to the ETR bracket of its parent group, not its own individual ETR.  able~III covers 12 ETR brackets (including two negative-profit categories) and reports entity counts and activity-share distributions within each bracket.

\paragraph{Key features}
Importantly, Table~III does \emph{not} contain financial aggregates (PROFIT, TAX\_ACCRUED, ASSETS, etc.), only entity counts (\code{ENTITIES\_COUNT}) and activity shares (\code{ACT\_*}). This makes it useful for understanding the composition of entities in each ETR range (e.g.\ the prevalence of IP-holding entities among very low-taxed MNE groups) but not for the direct revenue computation, which requires Table~IV.

\paragraph{Limitation}
The MNE-group ETR (used in Table~III) differs from the subsidiary-level ETR (used in Table~IV). As \citet{hugger2025} demonstrate, 53.2\% of all low-taxed profit (below 15\%) is located in jurisdictions where the average ETR exceeds 15\% at the MNE-group level. Table~III therefore significantly understates the share of entities facing low effective rates because it averages across the group's entire global portfolio.

\subsection{Table IV: Aggregate Totals by ETR Bracket (Subsidiary Level)}\label{sec:table4}

\paragraph{Purpose and structure}
Table~IV is the most policy-relevant table for GMT revenue simulation. It reports all financial aggregates disaggregated by 11 ETR brackets at the \emph{subsidiary} (TRSUB) level. The unit of analysis is the individual constituent entity, not the MNE group, so the ETR assigned to each row reflects the entity's own jurisdictional effective tax rate. This is the data dimension that \citet{hugger2025} exploit to demonstrate within-jurisdiction ETR heterogeneity.

\paragraph{Key features}
The \code{RATE\_0-15} bracket directly identifies all profit, assets, payroll and tax reported by entities facing an ETR below 15\%: the GMT top-up base. In 2022, Table~IV reports 49 parent jurisdictions with data, of which 45 have positive observations in the \code{RATE\_0-15} bracket. Table~IV has 27 MEASURE codes, the full set of financial and activity variables, matching Table~I, but with fewer affiliate jurisdictions (154) because the ETR disaggregation is only reported for affiliates with sufficient MNE groups to pass the OECD's confidentiality threshold.

\paragraph{Coverage by year.}
\begin{table}[H]
\centering
\caption{Year-by-year coverage of Table~IV (ETR Bracket Aggregates)}
\label{tab:t4yearcoverage}
\begin{tabular}{@{}lrrrrrr@{}}
\toprule
Year & Parent & Total  & Positive & Zero & NaN & Positive \\
     & jurisdictions & rows & obs & obs & obs & rate (\%) \\
\midrule
2016 & 25 &  3,005 &  2,937 &     68 &    0 & 97.7 \\
2017 & 33 &  3,936 &  3,727 &    209 &    0 & 94.7 \\
2018 & 37 &  4,677 &  4,363 &    314 &    0 & 93.3 \\
2019 & 41 & 21,604 & 17,377 &  4,227 &    0 & 80.4 \\
2020 & 43 & 20,096 & 16,393 &  3,703 &    0 & 81.6 \\
2021 & 42 & 26,861 & 16,404 & 13,893 & 6,564 & 61.1 \\
2022 & 45 & 21,220 & 17,154 &  4,058 &    8 & 80.8 \\
\midrule
\multicolumn{7}{l}{\small\textit{Note: The jump in total rows between 2018 and 2019 reflects the expansion}} \\
\multicolumn{7}{l}{\small\textit{of the affiliate jurisdiction dimension as more countries joined the}} \\
\multicolumn{7}{l}{\small\textit{OECD CbCR framework. The 2021 NaN spike mirrors Table~I (see Section~\ref{sec:table1}).}} \\
\bottomrule
\end{tabular}
\end{table}

A data quality issue affects Switzerland (CHE) in Table~IV for fiscal years 2017, 2020, and 2022. The reported profit values for Switzerland in these years are approximately 8,000-10,000 times larger than the corresponding Table~I aggregate totals. Specifically:

\begin{itemize}
    \item In 2022, CHE reports \code{PROFIT} in the \code{RATE\_0-15} bracket
          equal to USD~944 trillion (Table~IV), compared to a total worldwide
          profit of USD~109 billion in Table~I for 2022, a discrepancy of a
          factor of approximately 8,600.
    \item The anomaly is present in 2017 (factor $\approx$~10,000$\times$)
          and 2020 (factor $\approx$~9,000$\times$), but \emph{not} in
          2018 or 2021, where CHE Table~IV values are consistent with Table~I.
    \item The years 2018 and 2021 show CHE Table~IV profit in the 0-15\%
          bracket of approximately USD~60 billion and USD~92 billion
          respectively, both consistent with Switzerland's profile in Table~I
          as a medium-sized parent jurisdiction.
\end{itemize}

This anomaly likely reflects an error in the Swiss national CbCR submission for
those three years, rather than a genuine economic observation. The values exceed
total global GDP by orders of magnitude and are clearly erroneous. \textbf{For any
analysis using Table~IV, Switzerland should be excluded from ETR bracket
computations for fiscal years 2017, 2020, and 2022}, and its Table~I aggregate
totals should be used instead with an imputed ETR distribution.

\subsection{Table V: MNE Group Size Distribution}\label{sec:table5}

\paragraph{Purpose and structure}
Table~V reports the size distribution of MNE groups within each parent jurisdiction, measured by consolidated unrelated-party revenue (UPR), total assets, and employment. Unlike all other tables, Table~V has no \code{COUNTERPART\_AREA} dimension (it is set to \code{\_Z}): the size distribution is computed across all MNE groups headquartered in the parent jurisdiction, irrespective of where their affiliates are located.

\paragraph{Key features}
Table~V is the smallest table (4,195 rows) and has the highest positive observation rate (98.4\%), as percentile statistics are always defined where there are MNE groups. It provides seven percentile points (P5, P10, P25, P50, P75, P90, P95), enabling full characterisation of the within-jurisdiction MNE size distribution. The 40 parent jurisdictions covered is lower than in Table~I because the OECD requires a minimum number of MNE groups to compute meaningful percentiles.

\paragraph{Implications}
The UPR size distribution from Table~V provides the empirical basis for characterising the MNE groups at the borderline of the EUR~750~million threshold; those in the \code{SIZE\_LESS\_2\_5} class in Table~VI. For EU Member States, the P10 and P25 of UPR indicate how many MNE groups are within, say, two to three times the threshold, which determines the sensitivity of coverage to threshold changes.

\subsection{Table VI: Aggregate Totals by MNE Group Size Class}\label{sec:table6}

\paragraph{Purpose and structure}
Table~VI reports all financial and activity variables broken down by five MNE group revenue size classes (see Table~\ref{tab:mnesize}). It is the most recently added dimension of the CbCR database and currently covers fiscal year 2022 only.

\paragraph{Key features}
Table~VI covers 39 parent jurisdictions and 215 affiliate jurisdictions, with 27 MEASURE codes (the full set of financial and activity variables, minus \code{PROFIT\_ADJ} which is not available by size class). With 105,753 rows and a positive observation rate of 72.1\%, Table~VI provides reasonable coverage for the major MNE-hosting jurisdictions.

\paragraph{Use for methodology}
The size breakdown enables analysis of how the GMT's impact differs across the size distribution of MNE groups. In particular, it allows computation of SBIE ratios separately for large and medium MNE groups, testing the prediction of \citet{schjelderup2024} that the production subsidy effect is more pronounced among the largest, most capital-intensive groups.

\subsection{Reduced affiliate jurisdiction coverage in Table IV}

While Table~I covers 240 affiliate jurisdictions, Table~IV covers only 154. The
``missing'' 86 jurisdictions are primarily small or lower-income countries with
few MNE affiliates, for which the OECD applies confidentiality thresholds that
prevent disaggregation by ETR bracket. For the EU analysis in Option~A, this is
not a binding constraint: all 27 EU Member States appear in Table~IV as affiliate
jurisdictions for the relevant parent jurisdictions.

\subsection{Temporal lag}

The most recent CbCR release covers fiscal year 2022, published in
2025.\footnote{The 2025 OECD Corporate Tax Statistics release, published in
November 2025, is the source for the data described in this annex.} This implies a
lag of approximately three years between the fiscal year and public availability.
The 2023 data (which would cover the first year of GMT implementation) are not yet
available and, based on the historical pattern, are unlikely to be released before
mid-2026. The methodology is therefore necessarily backward-looking: it estimates
GMT exposure using pre-implementation data (2022) and applies the GMT rules
prospectively.



%----------------------------------------------------------------------
\section{OECD Statutory Corporate Income Tax Rates}\label{sec:cit}
%----------------------------------------------------------------------

The OECD \emph{Tax Database} provides harmonised time series of statutory
corporate income tax (CIT) rates, compiled from legislation and administrative
sources.  The statutory rate is the headline rate applied to corporate profits
before allowances and credits; it differs from the effective rates documented
in Section~\ref{sec:etr}.  Statutory rates are the primary instrument of tax
competition because they are the publicly visible signal to potential investors.

\subsection{Statutory CIT rates}

\textbf{File:} \texttt{raw/oecd\_tax\_rates/oecd\_statutory\_cit\_rates.csv}\\
\textbf{OECD dataflow:} \texttt{OECD.CTP.TPS,DSD\_TAX\_CIT@DF\_CIT,2.0}\\
\textbf{Rows:} 13{,}273 \quad \textbf{Size:} 1.2~MB\\
\textbf{Years:} 2000--2025 \quad \textbf{Countries:} 144

\begin{longtable}{@{}L{2.8cm} L{6.0cm} C{1.8cm} C{2.5cm}@{}}
\caption{Statutory CIT rates --- variable codes}\label{tab:cit_rates}\\
\toprule
\textbf{Code (MEASURE)} & \textbf{Description} & \textbf{Unit} & \textbf{Coverage} \\
\midrule
\endfirsthead
\toprule
\textbf{Code} & \textbf{Description} & \textbf{Unit} & \textbf{Coverage} \\
\midrule
\endhead
\bottomrule
\endlastfoot
\code{CIT}       & Combined (central + sub-central) statutory CIT rate & \% of income & 2000--2025 \\
\code{CIT\_C}    & Central government CIT rate only & \% of income & 2000--2025 \\
\code{CIT\_SCG}  & Sub-central government CIT rate (e.g.\ canton, state) & \% of income & 2000--2025 \\
\code{CIT\_EX\_ST} & Surtax or extra statutory charges & \% of income & 2000--2025 \\
\code{CIT\_LSN}  & Local surtax / net wealth tax equivalent included in CIT & \% of income & 2000--2025 \\
\end{longtable}

\noindent\textbf{Dimension notes:}
\begin{itemize}[noitemsep]
    \item \texttt{TARGETING}: \texttt{ST} = standard rate; \texttt{RED} = reduced rate; \texttt{INC} = increased rate.
    \item \texttt{SECTOR}: \texttt{S1311} = general government; \texttt{\_Z} = all.
    \item \texttt{DIV\_TREATMENT}: how intra-group dividends are treated.
\end{itemize}

The 144~countries include all OECD members plus a large number of non-OECD jurisdictions tracked by the OECD CTS programme, covering tax havens, Caribbean territories, and major emerging markets.  For analytical purposes the standard combined rate (\texttt{CIT}, \texttt{TARGETING=ST}) is typically the relevant measure.

\subsection{Combined statutory tax rates on dividend income}

\textbf{File:} \texttt{raw/oecd\_tax\_rates/oecd\_combined\_dividend\_rates.csv}\\
\textbf{OECD dataflow:} \texttt{OECD.CTP.TPS,DSD\_TAX\_CIT@DF\_CIT\_DIVD\_INCOME,1.0}\\
\textbf{Rows:} 7{,}915 \quad \textbf{Size:} 0.8~MB\\
\textbf{Years:} 2000--2025 \quad \textbf{Countries:} 38 (OECD members)

These data capture the combined burden of corporate, personal, and withholding taxes on dividend income, capturing the overall tax wedge on equity returns.

\begin{longtable}{@{}L{2.8cm} L{6.5cm} C{1.8cm} C{2.5cm}@{}}
\caption{Combined dividend tax rates --- variable codes}\label{tab:div_rates}\\
\toprule
\textbf{Code (MEASURE)} & \textbf{Description} & \textbf{Unit} & \textbf{Coverage} \\
\midrule
\endfirsthead
\toprule
\textbf{Code} & \textbf{Description} & \textbf{Unit} & \textbf{Coverage} \\
\midrule
\endhead
\bottomrule
\endlastfoot
\code{CITPITCIT}  & Integrated CIT + PIT rate on grossed-up dividend (classical system) & \% & 2000--2025 \\
\code{CPITCIT}    & Combined CIT + PIT rate (with integration relief) & \% & 2000--2025 \\
\code{CIT\_DP}    & CIT rate applied to dividend payment & \% & 2000--2025 \\
\code{DP}         & Dividend payment (per unit of pre-tax profit) & ratio & 2000--2025 \\
\code{FWHT}       & Final withholding tax rate on dividends & \% & 2000--2025 \\
\code{GUD}        & Grossed-up dividend (per unit of after-tax profit) & ratio & 2000--2025 \\
\code{IDTC}       & Imputation/dividend tax credit & \% & 2000--2025 \\
\code{IR}         & Imputation rate under partial imputation systems & \% & 2000--2025 \\
\code{NPT}        & Net personal tax on dividend income & \% & 2000--2025 \\
\code{PITGUD}     & Personal income tax on grossed-up dividend & \% & 2000--2025 \\
\code{PTDP}       & Personal tax on dividend payment & \% & 2000--2025 \\
\end{longtable}


%----------------------------------------------------------------------
\section{OECD Forward-Looking Effective Tax Rates}\label{sec:etr}
%----------------------------------------------------------------------

Forward-looking effective tax rates (ETRs) measure the tax burden on a
\emph{hypothetical marginal or average investment project}, incorporating not
only the statutory rate but also all relevant reliefs: capital allowances,
depreciation rules, interest deductibility, and targeted incentives such as
IP~boxes and R\&D tax provisions.  The OECD \emph{Corporate Tax Statistics ---
Forward-Looking Effective Tax Rates} dataset provides two complementary measures:

\begin{itemize}[noitemsep]
    \item \textbf{EMTR} (\emph{effective marginal tax rate}): the proportional tax
    wedge between the pre-tax and after-tax rate of return on the last unit of a
    break-even (zero-NPV) investment.
    \item \textbf{EATR} (\emph{effective average tax rate}): the ratio of the
    present value of taxes to the present value of pre-tax profit for a
    fixed-return investment.  Captures the margin relevant for location choice.
\end{itemize}

Three files are collected: a baseline (no incentives), a set incorporating
\emph{income-based} incentives (IP~boxes, patent boxes), and a set incorporating
\emph{expenditure-based} incentives (R\&D credits, enhanced allowances).

\subsection{ETR baseline}

\textbf{File:} \texttt{raw/oecd\_etr/oecd\_etr\_baseline.csv}\\
\textbf{OECD dataflow:} \texttt{OECD.CTP.TPS,DSD\_ETR@DF\_ETR\_BASELINE,2.0}\\
\textbf{Rows:} 29{,}489 \quad \textbf{Size:} 3.0~MB\\
\textbf{Years:} 2017--2024 \quad \textbf{Countries:} 103

\begin{longtable}{@{}L{3.2cm} L{6.0cm} C{2.0cm} C{2.5cm}@{}}
\caption{ETR baseline --- variable codes}\label{tab:etr_baseline}\\
\toprule
\textbf{Code (MEASURE)} & \textbf{Description} & \textbf{Unit} & \textbf{Coverage} \\
\midrule
\endfirsthead
\toprule
\textbf{Code} & \textbf{Description} & \textbf{Unit} & \textbf{Coverage} \\
\midrule
\endhead
\bottomrule
\endlastfoot
\code{EATR}               & Effective average tax rate (no incentives) & \% of income & 2017--2024 \\
\code{EMTR}               & Effective marginal tax rate (no incentives) & \% of income & 2017--2024 \\
\code{PBAR}               & $\bar{p}$: pre-tax rate of return on break-even investment & \% & 2017--2024 \\
\code{CAPITAL\_ALLOWANCES}& Present value of capital allowances (NPV of depreciation deductions) & \% of cost & 2017--2024 \\
\end{longtable}

\noindent\textbf{Dimension notes:}
\begin{itemize}[noitemsep]
    \item \texttt{ETR\_TAX\_BASIS}: \texttt{BASELINE} = no incentives; distinguishes \texttt{CORPORATE} from \texttt{SHAREHOLDER}-level.
    \item \texttt{ETR\_TAX\_TYPE}: \texttt{CS} = consolidated CIT; \texttt{FIXED} = fixed-rate assumption; \texttt{AWS} = actual weighted-average.
    \item \texttt{UNIT\_MEASURE}: \texttt{PT\_INC\_TAX} = as \% of income tax; \texttt{PT\_P5L} = as \% of initial cost.
\end{itemize}

\subsection{ETR with income-based incentives (IP~boxes)}

\textbf{File:} \texttt{raw/oecd\_etr/oecd\_etr\_income\_incentives.csv}\\
\textbf{OECD dataflow:} \texttt{OECD.CTP.TPS,DSD\_ETR@DF\_ETR\_IBRD,2.0}\\
\textbf{Rows:} 5{,}784 \quad \textbf{Size:} 0.7~MB\\
\textbf{Years:} 2000--2024 \quad \textbf{Countries:} 55 (with active IP regimes)

\begin{longtable}{@{}L{3.2cm} L{6.0cm} C{2.0cm} C{2.5cm}@{}}
\caption{ETR with income-based incentives --- variable codes}\label{tab:etr_income}\\
\toprule
\textbf{Code (MEASURE)} & \textbf{Description} & \textbf{Unit} & \textbf{Coverage} \\
\midrule
\endfirsthead
\toprule
\textbf{Code} & \textbf{Description} & \textbf{Unit} & \textbf{Coverage} \\
\midrule
\endhead
\bottomrule
\endlastfoot
\code{EATR}       & EATR \emph{without} income-based incentive applied & \% & 2000--2024 \\
\code{EATR\_ENH}  & EATR \emph{with} income-based incentive (e.g.\ IP~box rate) & \% & 2000--2024 \\
\code{PBAR}       & Pre-tax break-even rate \emph{without} incentive & \% & 2000--2024 \\
\code{PBAR\_ENH}  & Pre-tax break-even rate \emph{with} income-based incentive & \% & 2000--2024 \\
\end{longtable}

Each combination is reported per \texttt{REGIME}: \texttt{REGIME1} through \texttt{REGIME8} label specific IP~box or patent~box regimes in a given country-year; \texttt{\_N} indicates the standard (non-incentive) regime.

\subsection{ETR with expenditure-based incentives (R\&D)}

\textbf{File:} \texttt{raw/oecd\_etr/oecd\_etr\_expenditure\_incentives.csv}\\
\textbf{OECD dataflow:} \texttt{OECD.CTP.TPS,DSD\_ETR@DF\_ETR\_XBRD,2.0}\\
\textbf{Rows:} 1{,}224 \quad \textbf{Size:} 0.1~MB\\
\textbf{Years:} 2019--2024 \quad \textbf{Countries:} 51

Variable codes are identical to Table~\ref{tab:etr_income}, but the \texttt{ENH} variants incorporate expenditure-based R\&D tax credits or enhanced capital allowances rather than preferential income rates.  This dataset enables comparison of the incentive value of volume credits versus IP~box regimes.


%----------------------------------------------------------------------
\section{Wage and Earnings Data}\label{sec:wages}
%----------------------------------------------------------------------

Labour cost data are collected from three sources to maximise country coverage.
OECD average annual wages provide high-quality, internationally comparable data
for OECD member countries; the OECD hourly earnings index provides a longer
time series; and ILO hourly and monthly earnings extend coverage to a wider set
of developing countries.  These data serve as control variables in
profit-shifting regressions (wages proxy for labour market conditions) and as
inputs for constructing wage-to-profit ratios and SBIE payroll components.

\subsection{OECD average annual wages}

\textbf{File:} \texttt{raw/oecd\_wages/oecd\_avg\_annual\_wages.csv}\\
\textbf{OECD dataflow:} \texttt{OECD.ELS.SAE,DSD\_EARNINGS@AV\_AN\_WAGE,1.0}\\
\textbf{Rows:} 3{,}924 \quad \textbf{Size:} 0.4~MB\\
\textbf{Years:} 1990--2024 \quad \textbf{Countries:} 41 (OECD members + OECD aggregate)

\noindent\textbf{Key variable:} \code{WG} --- average annual wages of full-time dependent employees, denominated in national currency and in USD PPP (2024 constant).

\noindent\textbf{Dimension notes:}
\begin{itemize}[noitemsep]
    \item \texttt{UNIT\_MEASURE}: \texttt{USD\_PPP} (2024 constant PPP USD) is the primary cross-country comparator.
    \item \texttt{PRICE\_BASE}: \texttt{Q} = current prices; \texttt{R} = constant (real) prices.
    \item \texttt{SEX}: total (\texttt{\_Z}), male, or female.
\end{itemize}

\subsection{OECD hourly earnings index}

\textbf{File:} \texttt{raw/oecd\_wages/oecd\_hourly\_earnings.csv}\\
\textbf{OECD dataflow:} \texttt{OECD.SDD.TPS,DSD\_EAR@DF\_HOU\_EAR,1.0}\\
\textbf{Rows:} 41{,}421 \quad \textbf{Size:} 3.4~MB\\
\textbf{Years:} 1955--2026 \quad \textbf{Countries/groups:} 41

\noindent\textbf{Key variable:} \code{EAR} --- index of hourly earnings for employees (base year 2015 = 100).  This series is an \emph{index} rather than a level; for cross-country wage comparisons the OECD average annual wages dataset is more appropriate.  Coverage includes area aggregates: \texttt{EA19} (euro area), \texttt{EU27\_2020} (EU~27), \texttt{G7}, \texttt{OECD}.

\subsection{ILO global hourly earnings}

\textbf{File:} \texttt{raw/ilo\_wages/ilo\_hourly\_earnings.csv}\\
\textbf{ILO dataflow:} \texttt{ILO,DF\_EAR\_EHRA\_SEX\_ECO\_CUR\_NB,1.0}\\
\textbf{Rows:} 227{,}317 \quad \textbf{Size:} 50.4~MB\\
\textbf{Years:} 1990--2025 \quad \textbf{Countries:} 105

\noindent\textbf{Key variable:} \code{EAR\_EHRA\_NB} --- mean nominal hourly earnings of employees, denominated in local currency (\texttt{CUR\_TYPE\_LCU}) or PPP USD (\texttt{CUR\_TYPE\_PPP}).  Data are compiled from national household surveys, enterprise surveys, and administrative records reported to the ILO; methodological notes are stored in \texttt{NOTE\_SOURCE} and \texttt{NOTE\_INDICATOR} columns.

\noindent\textbf{Dimension notes:}
\begin{itemize}[noitemsep]
    \item \texttt{SEX}: \texttt{SEX\_T} = total; also male and female breakdowns.
    \item \texttt{ECO}: economic sector (\texttt{ECO\_SECTOR\_TOTAL} = economy-wide).
    \item \texttt{OBS\_STATUS}: \texttt{B} = break in series; \texttt{E} = estimated.
    \item Confidence bounds reported in \texttt{UPPER\_BOUND} / \texttt{LOWER\_BOUND} where available.
\end{itemize}

\subsection{ILO global monthly earnings}

\textbf{File:} \texttt{raw/ilo\_wages/ilo\_monthly\_earnings.csv}\\
\textbf{ILO dataflow:} \texttt{ILO,DF\_EAR\_EMTA\_SEX\_ECO\_CUR\_NB,1.0}\\
\textbf{Rows:} 529{,}234 \quad \textbf{Size:} 134~MB\\
\textbf{Years:} 1990--2025 \quad \textbf{Countries:} $\sim$140

\noindent\textbf{Key variable:} \code{EAR\_EMTA\_NB} --- mean nominal monthly earnings of employees.  Column structure is identical to the hourly earnings file.  Monthly earnings data exist for a wider set of countries than hourly earnings, making this file useful for non-OECD economies where hourly data are sparse.  The ILO \texttt{CUR\_TYPE\_PPP} observations provide PPP-adjusted values for cross-country analysis.

\noindent\textbf{Coverage note:} Coverage is uneven across countries and years; many non-OECD observations are spaced several years apart.  The OECD wage data should be preferred for OECD countries; ILO data fill the gap for non-OECD members.


%----------------------------------------------------------------------
\section{World Bank Macroeconomic Controls}\label{sec:worldbank}
%----------------------------------------------------------------------

World Bank data provide three types of country-level controls commonly included
in cross-country tax regressions: GDP per capita (economic development and
market size), population (country size), and income group classification (for
sub-group and heterogeneity analyses).

\subsection{GDP per capita}

\textbf{File:} \texttt{raw/world\_bank/wb\_gdp\_per\_capita.csv}\\
\textbf{Indicator:} \texttt{NY.GDP.PCAP.CD} (GDP per capita, current USD)\\
\textbf{Rows:} $\approx$2{,}660 \quad \textbf{Size:} $<$0.1~MB\\
\textbf{Years:} 2015--2024 \quad \textbf{Countries/groups:} $\approx$217

\noindent Columns: \texttt{iso3}, \texttt{country}, \texttt{year}, \texttt{gdp\_pc\_usd}.  The file includes World Bank regional and income-group aggregates (e.g.\ \texttt{ZH}: Africa Eastern and Southern), which should be filtered when merging with country-level data.

\subsection{Population}

\textbf{File:} \texttt{raw/world\_bank/wb\_population.csv}\\
\textbf{Indicator:} \texttt{SP.POP.TOTL} (total population)\\
\textbf{Rows:} $\approx$2{,}660 \quad \textbf{Size:} $<$0.1~MB\\
\textbf{Years:} 2015--2024 \quad \textbf{Countries/groups:} $\approx$217

\noindent Columns: \texttt{iso3}, \texttt{country}, \texttt{year}, \texttt{population}.

\subsection{Income group classifications}

\textbf{File:} \texttt{raw/world\_bank/wb\_income\_classifications.csv}\\
\textbf{Source:} World Bank country metadata API (\texttt{/v2/country})\\
\textbf{Rows:} 296 \quad \textbf{Type:} cross-section (current classification as of download)

\noindent Columns: \texttt{iso3}, \texttt{iso2}, \texttt{name}, \texttt{region}, \texttt{income\_level} (High / Upper-middle / Lower-middle / Low / Aggregates), \texttt{income\_id} (\texttt{HIC} / \texttt{UMC} / \texttt{LMC} / \texttt{LIC} / \texttt{NA}).

For analyses requiring time-varying income classifications, the \texttt{wb\_historical\_income\_groups.xlsx} file contains the full historical series from the OGHIST spreadsheet.


%----------------------------------------------------------------------
\section{OECD BEPS Implementation Indicators}\label{sec:beps}
%----------------------------------------------------------------------

The OECD BEPS project produced 15 Actions to curb aggressive tax planning.
The OECD's Questionnaire on Quality and Distribution of Data (QDD) programme
tracks countries' implementation of several key Actions: Action~2 (hybrid
mismatches), Action~4 (interest limitation), Action~5 (IP~regimes),
Action~12 (mandatory disclosure), and Action~13 (CbCR requirements), as well
as Controlled Foreign Company (CFC) rules.  These datasets provide
\emph{policy indicator} variables coding whether specific laws are in place,
their thresholds, and when they were introduced.

\noindent\textbf{Note on row counts:} Row counts below reflect the March~2026
download.  Several BEPS datasets have been materially expanded since earlier
versions (e.g.\ CFC rules grew from $\approx$1{,}400 to $\approx$1{,}900 rows
as additional years and countries were added); always verify against the actual
file when replicating.

\subsection{CbCR requirements (Action~13)}

\textbf{File:} \texttt{raw/pillar2\_beps/oecd\_cbcr\_requirements.csv}\\
\textbf{OECD dataflow:} \texttt{OECD.CTP.TPS,DSD\_QDD\_CBCREQ@DF\_QDD\_CBCREQ,1.0}\\
\textbf{Rows:} 493 \quad \textbf{Countries:} 119

Key variables: \code{CBCR\_LAW} (whether country has enacted domestic CbCR law, Yes/No), \code{THRESHOLD} (revenue threshold, EUR millions), \code{DEADLINE} (filing deadline, months after fiscal year-end), \code{UPE} (ultimate parent entity jurisdiction indicator).

\subsection{CFC rules}

\textbf{File:} \texttt{raw/pillar2\_beps/oecd\_cfc\_rules.csv}\\
\textbf{OECD dataflow:} \texttt{OECD.CTP.TPS,DSD\_QDD\_CFC@DF\_QDD\_CFC,2.0}\\
\textbf{Rows:} 1{,}901 \quad \textbf{Years:} 2019, 2024, 2025 \quad \textbf{Countries:} 146

CFC rules tax the passive income of foreign subsidiaries at the parent level when the subsidiary pays little or no tax.  Key variables: \code{CFC} (Yes/No), \code{TRIGGER\_RATE} (low-tax threshold rate triggering CFC charge, \%), \code{YEAR} (year of first introduction), \code{SUB} (substance exclusions).

\subsection{Hybrid mismatch rules (Action~2)}

\textbf{File:} \texttt{raw/pillar2\_beps/oecd\_hybrid\_mismatch.csv}\\
\textbf{OECD dataflow:} \texttt{OECD.CTP.TPS,DSD\_QDD\_ACT2@DF\_QDD\_ACT2,1.0}\\
\textbf{Rows:} 1{,}171 \quad \textbf{Year:} 2025 \quad \textbf{Countries:} 145

Key variables: \code{ACT2} (whether Action~2 rules are in place), \code{BR\_REC1}--\code{BR\_REC5} (implementation of each primary branch rule), \code{REC1}--\code{REC15} (implementation of specific OECD recommendations), \code{YEAR} (year of legislative enactment).

\subsection{Interest limitation rules (Action~4)}

\textbf{File:} \texttt{raw/pillar2\_beps/oecd\_interest\_limitation.csv}\\
\textbf{OECD dataflow:} \texttt{OECD.CTP.TPS,DSD\_QDD\_ILR@DF\_QDD\_ILR,2.0}\\
\textbf{Rows:} 3{,}879 \quad \textbf{Years:} 2019, 2024, 2025 \quad \textbf{Countries:} 146

Interest limitation rules cap deductible interest expense, typically to a fraction of EBITDA, curtailing debt-shifting.  Key variables: \code{ILR} (Yes/No), \code{GR} (general ratio threshold, e.g.\ 30\% of EBITDA), \code{TR} (safe harbour threshold, EUR), \code{FAM} (group ratio/equity escape, Yes/No), \code{CF} (carry-forward of disallowed interest, Yes/No), \code{CF\_YEARS} (maximum carry-forward years).

\subsection{Intellectual property regimes (Action~5)}

\textbf{File:} \texttt{raw/pillar2\_beps/oecd\_ip\_regimes.csv}\\
\textbf{OECD dataflow:} \texttt{OECD.CTP.TPS,DSD\_QDD\_IPR@DF\_QDD\_IPR,2.1}\\
\textbf{Rows:} 2{,}305 \quad \textbf{Years:} 2018--2025 \quad \textbf{Countries:} 64

IP~box (patent box) regimes tax qualifying IP income at a preferential rate.  BEPS Action~5 introduced the \emph{modified nexus approach}, linking relief to genuine R\&D expenditure.  Each regime is identified by \code{REGIME} (REGIME1, REGIME2, \ldots).  Key variables: \code{NAME} (regime name), \code{RATE} (preferential rate, \%), \code{ELI\_IP} (eligible IP assets), \code{NEXUS} (whether modified nexus approach adopted), \code{YEAR\_START}, \code{YEAR\_END}.

\subsection{Mandatory disclosure rules (Action~12)}

\textbf{File:} \texttt{raw/pillar2\_beps/oecd\_mandatory\_disclosure.csv}\\
\textbf{OECD dataflow:} \texttt{OECD.CTP.TPS,DSD\_QDD\_ACT12@DF\_QDD\_ACT12,1.0}\\
\textbf{Rows:} 711 \quad \textbf{Year:} 2025 \quad \textbf{Countries:} 145

Mandatory disclosure rules require intermediaries or taxpayers to proactively disclose potentially aggressive tax arrangements.  Key variables: \code{ACT12} (Yes/No), \code{RULE\_DESC} (description), \code{YEAR} (year of introduction).


%----------------------------------------------------------------------
\section{Pillar Two Implementation Status}\label{sec:pillar2}
%----------------------------------------------------------------------

BEPS Pillar Two introduces a global minimum effective tax rate of 15\% for
large MNE groups (consolidated revenue $\geq$ EUR~750~million) through
the Income Inclusion Rule (IIR), the Undertaxed Payments Rule (UTPR), and
the Qualified Domestic Minimum Top-up Tax (QDMTT).

\subsection{PwC Pillar Two country tracker}

\textbf{File:} \texttt{raw/pillar2\_qdmtt/pwc\_pillar2\_tracker.csv}\\
\textbf{Source:} PricewaterhouseCoopers global Pillar Two tracker\\
\textbf{Rows:} 1{,}190 \quad \textbf{Type:} panel (requirement $\times$ effective-date level)

Key columns: \texttt{Territory}, \texttt{Requirement} (\texttt{GIR}, \texttt{IIR}, \texttt{UTPR}, \texttt{QDMTT}), \texttt{P2 Start} (effective start date), \texttt{Status} (\texttt{Enacted}, \texttt{Consultation}, etc.), \texttt{Caveat} (qualitative notes on transition rules).

\noindent\textbf{Coverage note:} Only countries that have taken legislative steps appear in the tracker; non-adopters are absent.  Cross-referencing against the OECD Central Record of Qualified Domestic Minimum Taxes (PDF stored at \texttt{raw/pillar2\_qdmtt/oecd\_pillar2\_central\_record.pdf}) is advisable for publications.

\subsection{OECD central record (PDF)}

\textbf{File:} \texttt{raw/pillar2\_qdmtt/oecd\_pillar2\_central\_record.pdf}\\
\textbf{Source:} OECD Central Record of Qualified Domestic Minimum Taxes (transitional qualified status)

The PDF document records the official OECD assessment of which domestic minimum taxes and top-up tax rules qualify under the GloBE framework, providing an authoritative cross-check on the PwC tracker.


%----------------------------------------------------------------------
\section{OECD Comparative Tax Revenue Statistics}\label{sec:cts}
%----------------------------------------------------------------------

\textbf{File:} \texttt{raw/oecd\_cts\_additional/oecd\_cit\_revenue\_gdp.csv}\\
\textbf{OECD dataflow:} \texttt{OECD.CTP.TPS,DSD\_REV\_COMP\_GLOBAL@DF\_RSGLOBAL,2.1}\\
\textbf{Rows:} 3{,}217{,}432 \quad \textbf{Size:} 342~MB\\
\textbf{Years:} 1990--2024 \quad \textbf{Countries:} 143

This is the OECD \emph{Revenue Statistics --- Comparative tables} dataset, providing harmonised time series of tax revenue by type of tax, level of government, and country.  It is the primary source for CIT revenue as a share of GDP, total tax, and national accounts.

\noindent\textbf{Column structure:} \texttt{DATAFLOW}, \texttt{REF\_AREA} (ISO 3-letter), \texttt{MEASURE} (\code{TAX\_REV} for all rows), \texttt{SECTOR} (level of government: \texttt{S13}~=~general government, \texttt{S1311}~=~central, \texttt{S1312}~=~state/regional, \texttt{S1313}~=~local, \texttt{S1314}~=~social security), \texttt{STANDARD\_REVENUE} (OECD revenue classification code, e.g.\ \texttt{T\_2000} = corporate income taxes, \texttt{T\_1000} = total income taxes), \texttt{CTRY\_SPECIFIC\_REVENUE} (country-specific sub-classifications), \texttt{UNIT\_MEASURE} (\texttt{PT\_OTR\_SECTOR} = \% of GDP; \texttt{PC\_TAXREV} = \% of total tax revenue; \texttt{USD\_PPP} = PPP USD billions), \texttt{TIME\_PERIOD}, \texttt{OBS\_VALUE}, \texttt{REVENUE\_CODE}.

\noindent\textbf{Use in methodology:} The \texttt{T\_2000} series (\% of GDP) provides the empirical baseline for assessing how much of current CIT revenue is at risk from GMT implementation and how GMT top-up revenues compare with existing CIT receipts.  The \texttt{S13}/\texttt{S1311} distinction matters for federal states where corporate taxes are levied at both central and sub-national levels.

\noindent\textbf{File size note:} At 342~MB and 3.2~million rows, this is the largest single file in the collection.  The large size reflects the cross-product of 143 countries $\times$ 35 years $\times$ multiple revenue codes $\times$ government levels.  For most analyses, filtering to \texttt{UNIT\_MEASURE = PT\_OTR\_SECTOR} and \texttt{STANDARD\_REVENUE IN (T\_2000, T\_0000)} reduces the working set to manageable size.


%----------------------------------------------------------------------
\section{BEA US Multinational Enterprise Data}\label{sec:bea}
%----------------------------------------------------------------------

\textbf{Files:} \texttt{raw/bea\_mne/bea\_mofa\_*.xls} (four Excel workbooks)\\
\textbf{Source:} US Bureau of Economic Analysis (BEA), Survey of US Direct Investment Abroad\\
\textbf{Status:} Downloaded; four Excel workbooks present (March~2026)

The BEA provides detailed data on US-headquartered majority-owned foreign affiliates (MOFAs), including employment, assets, sales, and value added by country and industry.  These data complement the OECD CbCR statistics by providing a longer time series (back to the 1980s) for the US parent population, as well as industry-level decompositions not available in the CbCR.

\begin{table}[H]
\centering
\caption{BEA MOFA files downloaded}\label{tab:bea_files}
\begin{tabular}{@{}L{5.5cm} L{6.5cm} C{1.5cm}@{}}
\toprule
\textbf{File} & \textbf{Content} & \textbf{Size} \\
\midrule
\texttt{bea\_mofa\_selected\_by\_country.xls}  & Selected financial and activity data by country & 0.6~MB \\
\texttt{bea\_mofa\_employment.xls}             & Employment by country and industry & 0.5~MB \\
\texttt{bea\_mofa\_assets.xls}                 & Total assets by country and industry & 0.3~MB \\
\texttt{bea\_mofa\_value\_added.xls}           & Value added by country and industry & 0.4~MB \\
\bottomrule
\end{tabular}
\end{table}

\noindent\textbf{Limitations:} BEA MOFA data cover only US-parented MNEs.  They are not directly comparable to the OECD CbCR aggregates, which cover all reporting parent jurisdictions.  For EU-focused GMT analysis, the BEA data are most useful as an external validation of profit and employment levels in specific host countries.  An optional API path is also coded in \texttt{collect\_all\_data.py} (Series~IDs 5, 6, 8, 22, 25); this requires a free \texttt{BEA\_API\_KEY}.


%----------------------------------------------------------------------
\section{Eurostat Inward Foreign Affiliate Statistics (FATS)}\label{sec:fats}
%----------------------------------------------------------------------

\textbf{Canonical file:} \texttt{raw/eurostat\_fats/fats\_inward\_eu\_full.csv}\\
\textbf{Supporting raw files:} \texttt{raw/eurostat\_fats/fats\_g1b\_08.csv};\quad
\texttt{raw/eurostat\_fats/fats\_activ.csv}\\
\textbf{Eurostat datasets:} \texttt{fats\_g1b\_08} (2008--2020) and \texttt{fats\_activ} (2021 onward)\\
\textbf{Combined rows:} 4{,}623{,}851 \quad \textbf{Combined size:} 442~MB\\
\textbf{Years:} 2008--2023 \quad \textbf{Downloaded:} March~2026 from the live Eurostat API

Eurostat inward FATS provide information on foreign-controlled enterprises
operating inside European reporting economies.  The data are useful for
descriptive validation, sectoral context, and external checks on the location
of real activity.  They are not a substitute for the OECD CbCR data, because
they do not contain tax liabilities or effective tax rates.

\subsection{Source files and rebuild procedure}

The canonical file required a documented rebuild.  The repository initially
contained only \texttt{fats\_g1b\_08.csv}, which ended in 2020.  The live
Eurostat metadata showed that the public inward FATS series is split into two
official datasets: the legacy series \texttt{fats\_g1b\_08}
(\emph{Foreign control of enterprises by controlling countries, 2008--2020})
and the newer series \texttt{fats\_activ}
(\emph{Foreign controlled EU enterprises by NACE Rev.2 activity and country of
control, from 2021 onwards}).

The collection script \texttt{collect\_all\_data.py} was updated to download
both datasets through the Eurostat SDMX~2.1 API, preserve both raw files, and
build one canonical project file by vertically appending the two
non-overlapping year ranges.  Provenance columns were added to the canonical
file: \texttt{source\_dataset}, \texttt{indicator\_scheme}, and
\texttt{indicator\_code}, so the series break remains visible in any later
analysis.

\begin{table}[H]
\centering
\caption{Eurostat inward FATS raw files}\label{tab:fats_files}
\begin{tabular}{@{}L{5.5cm} C{1.8cm} C{2.0cm} C{2.0cm}@{}}
\toprule
\textbf{File} & \textbf{Rows} & \textbf{Years} & \textbf{Size} \\
\midrule
\texttt{fats\_g1b\_08.csv}          & 1{,}421{,}600 & 2008--2020 & 109~MB \\
\texttt{fats\_activ.csv}            & 3{,}202{,}251 & 2021--2023 & 241~MB \\
\texttt{fats\_inward\_eu\_full.csv} & 4{,}623{,}851 & 2008--2023 & 442~MB \\
\bottomrule
\end{tabular}
\end{table}

\subsection{Series break at 2021}

The merge was intentionally conservative.  The two source files do not overlap
in years, so the merge is a clean stacked append.  However, the two halves are
\emph{not} directly comparable across the 2021 break:

\begin{itemize}[noitemsep]
    \item \textbf{Activity dimension:} the old file has one broad aggregate
    (\texttt{B-N\_S95\_X\_K}), while the new file has 48 NACE-based activity
    categories plus a different top aggregate (\texttt{B-S\_X\_O\_S94}).
    \item \textbf{Indicator system:} the old file uses 23 legacy
    \texttt{indic\_sb} codes (e.g.\ \texttt{V12110}, \texttt{V12150}); the new
    file uses 19 \texttt{indic\_sbs} codes (e.g.\ \texttt{NETTUR\_MEUR},
    \texttt{AV\_MEUR}).
    \item \textbf{Country-of-control:} the old file has 255 control-country
    categories; the new file has 50.
    \item \textbf{Concept differences:} turnover becomes \emph{net} turnover;
    personnel costs become \emph{employee benefits expense}; the investment
    concept is narrowed to tangible non-current assets.
\end{itemize}

\noindent\textbf{Practical implications:} Eurostat FATS are best treated as a
supporting EU dataset.  Any figure spanning both sides of the 2021 break must
either restrict to clearly harmonised concepts or mark the break explicitly.  A
pre/post econometric design requires a separate harmonisation step mapping only
clearly comparable indicators across the break.  For the GMT thesis, the safest
use is to draw on the post-2021 \texttt{fats\_activ} structure for recent EU
sectoral evidence and use the longer combined file mainly for descriptive
background.


\newpage
\printbibliography % chktex 1

\end{document}
